\chapter{Ce qui fait encore mal — et le remède}
\label{ch:douleurs-restantes}

\section{Deux lectures du même mot}

Le mot \emph{douleur} désigne ici deux choses qu'il faut distinguer :

\begin{itemize}
  \item la douleur \textbf{résolue} (\S\ref{ch:douleurs}) — ce que l'assureur et
    le promoteur subissaient avant, et que l'architecture traite par
    construction ;
  \item la douleur \textbf{restante} — ce qui, aujourd'hui, continue de faire mal
    à l'usage, et qui n'est pas caché dans ce rapport.
\end{itemize}

Cette seconde liste est la plus utile au lecteur : elle dit où le produit
n'est pas encore au niveau de sa promesse, et ce qu'il faudrait décider pour y
arriver.

\section{Les huit douleurs restantes}

\begin{table}[h]
\centering
\small
\begin{tabularx}{\textwidth}{@{}p{4.3cm}X@{}}
\toprule
\textbf{Douleur} & \textbf{Effet pour l'assureur ou le promoteur} \\
\midrule
P1 Prose refusée selon le format des nombres & Le résumé exécutif peut retomber
  en gabarit tandis que le reste est rédigé : document inégal devant un client \\
P2 Prose non reproductible & Deux demandes identiques peuvent produire deux
  formulations différentes \\
P3 Tier gratuit peu compatible avec des adresses réelles & Le contenu envoyé
  peut servir à l'amélioration des produits du fournisseur \\
P4 Communication publique en décalage & La page publique crédite encore un autre
  fournisseur et promet une souveraineté que le tier gratuit n'offre pas \\
P5 Aléa argile communal & Impossible de distinguer deux bâtiments d'une même
  commune — précisément sur l'aléa le plus sinistrogène \\
P6 Cache non persistant en serverless & Chaque rapport peut consommer un appel
  de modèle : un pic peut épuiser un quota et basculer en gabarit en pleine
  démonstration \\
P7 Dette de nommage & Le point d'entrée garde un nom de fournisseur qui n'est
  plus celui par défaut \\
P8 Outillage de test fragile & La suite ne démarre pas sans exclure un plugin
  tiers bloqué par la politique de sécurité de la machine \\
\bottomrule
\end{tabularx}
\caption{Les huit douleurs restantes, du point de vue de l'utilisateur.}
\end{table}

\section{Les trois douleurs qui coûtent le plus}

\subsection{P1 — La prose dépend du format des nombres}

C'est la douleur la plus visible, et elle a été \textbf{mesurée}. Le validateur
d'ancrage n'accepte que les nombres présents dans le rapport. Or le fournisseur
actuel écrit volontiers les nombres sous forme compacte ou à la française
(« 409 k€ », « 1,1 m »), là où l'ensemble autorisé ne contient que les formes
complètes (\code{408700}, \code{294264}, \code{523136}).

Résultat concret : les cinq narratifs de catégorie et les quatre narratifs de
mitigation sont acceptés, mais le \emph{résumé exécutif} est refusé et remplacé
par le gabarit calculé.

\begin{douleur}[title={Douleur observée}]
Le document mélange une prose de qualité et un premier paragraphe mécanique. Le
client ne connaît pas la cause et la lit comme un défaut de finition — alors
qu'il s'agit d'un garde-fou qui fonctionne.
\end{douleur}

\begin{remede}[title={Remède proposé}]
Élargir l'ensemble des nombres autorisés aux \textbf{faits du rapport lui-même} :
profondeur du scénario, bornes basses et hautes des dommages, et leurs formes
compactes (k€), plus les chiffres portés par le libellé du scénario. Le garde-fou
n'est pas affaibli : on autorise des faits du document, jamais des chiffres
extérieurs. Le changement exige un test sur une sortie réelle du modèle, car
cette autorité d'ancrage est partagée avec l'assistant documentaire.
\end{remede}

\subsection{P3 — Un tier gratuit qui n'est pas compatible avec des adresses}

Un assureur manipule des adresses de biens : la donnée est identifiante. Or, sur
le tier gratuit du fournisseur retenu par défaut, le contenu envoyé
\emph{peut servir à l'amélioration des produits} du fournisseur. Le même
connecteur, avec une clé facturée, ne présente pas cette réserve.

\begin{douleur}[title={Douleur}]
Le service fonctionne, mais pas dans les conditions qu'un client assurance
attendrait pour des données de portefeuille.
\end{douleur}

\begin{remede}[title={Remède proposé}]
Le fournisseur est déjà réglable par configuration : basculer sur une clé
facturée ne demande aucune modification de code. La décision à prendre est
commerciale (qui paie l'appel), pas technique. À défaut, un fournisseur
souverain peut être branché sans toucher au reste de l'architecture, puisque
l'interface du connecteur est neutralisée par un simple routage.
\end{remede}

\subsection{P5 — L'aléa argile reste communal}

C'est la douleur la plus coûteuse \emph{en sinistralité}. Le retrait-gonflement
des argiles est aujourd'hui publié en \code{commune-level-estimate} : deux
bâtiments d'une même commune reçoivent la même estimation, alors que leur
exposition réelle dépend du sol, de la végétation et du bâti.

\begin{douleur}[title={Douleur}]
Sur l'aléa qui pèse le plus dans les sinistres de sécheresse, le service ne
différencie pas les biens — ce qui est exactement ce qu'on lui demande par
ailleurs.
\end{douleur}

\begin{remede}[title={Remède proposé}]
Câbler une source vectorielle d'aléa argile et passer la résolution à
\code{building-level} ou \code{point-in-polygon}. La règle actuelle est inscrite
dans la constitution du projet : la modification doit donc passer par un
amendement explicite, puis par la spécification — jamais par une retouche de
code silencieuse.
\end{remede}

\section{Les cinq autres, en bref}

\begin{table}[h]
\centering
\small
\begin{tabularx}{\textwidth}{@{}p{1.1cm}X@{}}
\toprule
\textbf{} & \textbf{Douleur \(\rightarrow\) remède} \\
\midrule
P2 & Prose non reproductible \(\rightarrow\) le cache rend une même demande
  identique au rejeu ; pour le reste, la reproductibilité n'est pas promise, et
  aucun fait ne dépend de la prose. \\
P4 & Communication publique en décalage \(\rightarrow\) mettre la page publique,
  la FAQ et la mention légale en accord avec le fournisseur réellement utilisé,
  ou revenir à un fournisseur souverain. \\
P6 & Cache non persistant en serverless \(\rightarrow\) stockage externe, ou
  budget d'appels de prose journalier, sur le modèle du budget BDNB déjà en
  place. \\
P7 & Dette de nommage du point d'entrée \(\rightarrow\) renommage neutre dans un
  change set dédié, qui met à jour les tests qui épinglent ce chemin d'appel. \\
P8 & Outillage de test fragile \(\rightarrow\) exclusion documentée du plugin
  tiers en échec, ou image d'intégration continue maîtrisée, pour que
  « lancer les tests » redevienne une commande unique. \\
\bottomrule
\end{tabularx}
\caption{Les cinq douleurs restantes et leur remède.}
\end{table}

\section{Les décisions qui reviennent au commanditaire}

Trois questions ne sont pas techniques et ne peuvent pas être tranchées par
l'équipe :

\begin{enumerate}
  \item \textbf{Qui paie l'appel au modèle ?} Le tier gratuit est utilisable et
    suffisant en démonstration ; il n'est pas compatible avec des adresses de
    portefeuille réelles.
  \item \textbf{Où doit vivre la prose ?} En Europe (souveraineté) ou chez le
    fournisseur le moins cher (coût) — l'architecture supporte les deux, la
    décision est commerciale.
  \item \textbf{Jusqu'où automatiser l'interprétation ?} Le projet s'arrête
    volontairement avant le score et la recommandation. Aller plus loin exige un
    amendement de la constitution, pas un ajout de champ.
  \end{enumerate}
